\section{Desafios e Problemas Encontrados}

O detector trabalha por \textit{timeout} e não distingue nó lento de nó morto. Quando um
\texttt{REQUEST} falha ou expira, o nó trata o silêncio como permissão: é o que impede o
\textit{cluster} de travar quando um \textit{peer} morre, mas o mesmo caminho dispara quando ele
apenas sofre uma pausa de GC. Reproduzimos o caso com dois nós reais, e ambos entraram na seção
crítica. Não corrigimos: separar lento de morto exige quórum ou consenso, e privilegiamos a
disponibilidade.

Já a queda de um \textit{peer} entre o \texttt{REQUEST} e o \texttt{REPLY} deixava a aquisição
presa, pois o nó aguarda $N-1$ permissões; passamos a reenviar o pedido a cada 500 ms e a
recontatar quem volta após a quarentena, o que também destravou a queda de seguidor, que custava
30 s por rodada. Configurações divergentes, por fim, geravam \textit{split-brain} silencioso e
agora são recusadas na subida, com código 2.
